\chapter{F1.a.2 Class-Match at Wild Primes and Paskunas Block-Centraliser
Regularity}
\label{v4-arith:w6-f:chapter}

\emph{At the wild primes \(p=2,3\), the class-level comparison is a
comparison between automorphic block \(\mathrm{Ext}^{1}\)-classes and
the tangent fibre of the wild-enhanced Emerton--Gee stack. The
non-semisimple Bushnell--Henniart blocks determine fixed-determinant
local deformation functors; their Zariski tangent spaces are the
corresponding \(H^{1}(G_{\mathbb{Q}_{p}},\mathrm{ad}^{0})\). Under the
wild-EG enhancement, the residual gerbe identifies this tangent space
with the degree-one derived tangent fibre, and the quiver arrows
\(\alpha_{01},\alpha_{12}\) are sent to the corresponding tangent
directions. The very-special ordinary multiplicities and the wild
block-centraliser regularity enter only as explicit conditional
inputs.}

\section{Objects}
\label{v4-arith:w6-f:domain}

\subsection{The Three Boundary Conditions}

The Paskunas comparison has three named boundary conditions:
\begin{itemize}
\item \textbf{F1.a.2.class-match}
\Cref{v4-arith:f1a3-BH:conj:Paskunas-wild-ext}): the class-level
identification (not merely dimensional) between the
\(\mathrm{Ext}^{1}\)-classes \(\alpha_{01}\) at \(p=2\) and
\((\alpha_{01},\alpha_{12})\) at \(p=3\) on the automorphic side and the
degree-\(1\) tangent fibre of the wild-enhanced Emerton--Gee stack.
\item \textbf{R.B.1}
\Cref{v4-arith:w5-f:F1a3-AB:prop:remainder-AB}(R.B.1)): the very-special
ordinary-block Breuil--Mezard multiplicity at \(p\geq 5\) on the one
sub-stratum where Paskunas 2015 did not give the explicit number.
\item \textbf{R.B.2}
\Cref{v4-arith:w5-f:F1a3-AB:prop:remainder-AB}(R.B.2)): the wild-prime
block centraliser regularity result, the analogue of Paskunas' \(p\geq 5\)
block-centraliser Cohen--Macaulay / regular local dimension result at
\(p\in\{2,3\}\).
\end{itemize}

\subsection{Inherited Objects}

All categorical inputs are prior:
\begin{itemize}
\item the block decomposition of \(\mathrm{Mod}^{\mathrm{sm},\chi}_{GL_{2}(\mathbb{Q}_{p})}\) at
\(p=2\) (\Cref{v4-arith:w5-e:thm:Paskunas-p2}) and at \(p=3\)
(\Cref{v4-arith:w5-e:thm:CDP-p3});
\item the BH type list
\(\mathbf{T}^{\mathrm{BH}}_{p,\chi}\)
(\Cref{v4-arith:f1a3-BH:prop:BH-types-p2}--\ref{v4-arith:f1a3-BH:prop:BH-types-p3});
\item the wild-enhanced Emerton--Gee stack
\(\mathcal{X}^{\mathrm{EG,wild},p,\chi}_{GL_{2}}\)
(\Cref{v4-arith:deg-f1a:def:EG-wild-enhanced});
\item the block-algebra level \(\mathrm{Ext}^{1}\) computations at
\(p=2\) (\Cref{v4-arith:bzn-rcompletion:prop:R4-Ext1-p2}) and
\(p=3\) (\Cref{v4-arith:bzn-rcompletion:prop:R4-Ext1-p3});
\item the Paskunas generic-block regularity statement
\Cref{v4-arith:w5-f:F1a3-AB:thm:Paskunas-centraliser}(1) at \(p\geq 5\).
\end{itemize}

\subsection{Additional Inputs}

The following inputs are logically separate from the preceding
automorphic block computations.
\begin{description}
\item[Very-special multiplicity datum.] The missing Breuil--Mezard
values in the very-special ordinary block at \(p\geq5\).
\item[Wild block-centraliser datum.] The Cohen--Macaulay and generic
regularity statement for the wild blocks at \(p=2,3\).
\item[Colmez--Dospinescu--Paskunas 2014, \S 5.]
\(p\)-adic local Langlands at \(p=3\); contains the block centraliser
structure for the three principal-block types.
\item[Bella\"iche--Chenevier 2009.] \emph{Families of Galois
representations and Selmer groups}, Ast\'erisque 324. The
residual-gerbe Galois deformation functor at a generic point.
\item[Kisin 2008 J.~AMS~21.] ``Potentially semi-stable
deformation rings''. The tangent and cotangent-complex structure of local
Galois deformation rings at \(GL_{2}(\mathbb{Q}_{p})\).
\item[Emerton--Gee arXiv:1908.07185 \S 9.] Wild inertia block
structure on the spectral side; the component-level description of
the wild-enhanced EG stack at \(p\in\{2,3\}\) with explicit tangent
sheaves.
\end{description}

\subsection{Independence of Sources}

The earlier proof sources for F1.a.2 and F1.a.3.A+B are:
Pa{\v s}k{\=u}nas 2013, Colmez--Dospinescu--Pa{\v s}k{\=u}nas 2014,
Bushnell--Henniart 2006, Pa\v skunas 2015, Emerton--Pa\v skunas 2015,
Jantzen 2003, and Carter 1989.
The present conditional analysis draws on: the two data above,
Bella\"iche--Chenevier 2009, Kisin 2008, and Emerton--Gee
arXiv:1908.07185 \S 9.

The very-special datum concerns the sub-stratum for which Pa\v skunas
2015 does not give the multiplicity. The wild centraliser datum supplies the
Cohen--Macaulay and generic-regularity information not proved in
the preceding block computations.
Bella\"iche--Chenevier 2009 is the pre-EG residual-gerbe Galois deformation framework,
independent of Pa\v skunas' block method.
Kisin 2008 is local deformation theory for \(GL_{2}(\mathbb{Q}_{p})\), independent of the
automorphic block construction.
Emerton--Gee \S 9 is the wild-inertia component analysis on the spectral side, disjoint
from the automorphic block method.

\textit{Note on Colmez--Dospinescu--Pa\v skunas 2014 \S 5.}
This paper appears in two distinct roles: its \S 5.3 connecting-map result (the
Colmez--Montreal-functor identification of block-Ext with Galois-Ext) serves as the
automorphic input for the class-match theorem
(\Cref{v4-arith:w6-f:prop:Paskunas-equals-Galois}). A separate
centraliser input would be needed for R.B.2
(\Cref{v4-arith:w6-f:thm:CDP-p3-centraliser}).
These two uses draw on distinct propositions within \S 5.3. The
connecting-map result and the centraliser-dimension results are
mathematically independent within the paper.

\section{F1.a.2 Class-Match at Wild Primes}
\label{v4-arith:w6-f:class-match}

\subsection{The Class-Level Target}

The class-level compatibility statement is the following: each
non-semisimple block's \(\mathrm{Ext}^{1}\)-class on the automorphic side
must correspond, as a specific \(k\)-line, to a specific degree-\(1\)
tangent direction on the wild-enhanced EG stack.
We write
\[
\mathbb{T}^{1}_{\mathcal X}|_{x}:=
H^{0}\bigl((\mathbb L_{\mathcal X}|_{x})^{\vee}\bigr)
\]
for this infinitesimal tangent fibre.

\begin{definition}[the class-match target, \(p=2\)]
\label{v4-arith:w6-f:def:class-match-p2}
\ClaimStatusDefinitional. At \(p=2\), fix \(\chi\) and the unique
non-semisimple principal block \(\mathfrak{B}_{2}\). The
\emph{class-match target} is a \(k\)-linear isomorphism
\begin{equation}
\label{v4-arith:w6-f:eq:class-match-p2}
\mathsf{c}^{\mathrm{class}}_{2}\colon
\mathrm{Ext}^{1}_{\mathfrak{B}_{2}}(\mathbf{1}_{\mathrm{triv}},\mathbf{1}_{\mathrm{sgn}})
\;\xrightarrow{\;\sim\;}\;
\mathbb{T}^{1}_{\mathcal{X}^{\mathrm{EG,wild},2,\chi}_{GL_{2}}}\big|_{x_{\tau_{2}}},
\end{equation}
sending the Paskunas \(\mathrm{Ext}^{1}\)-generator \(\alpha\) of
\Cref{v4-arith:bzn-rcompletion:prop:R4-Ext1-p2} to the unique
tangent direction at the reducible closed point \(x_{\tau_{2}}\).
\end{definition}

\begin{definition}[the class-match target, \(p=3\)]
\label{v4-arith:w6-f:def:class-match-p3}
\ClaimStatusDefinitional. At \(p=3\), fix \(\chi\) and the three
principal-block types \(\tau_{3,0},\tau_{3,1},\tau_{3,2}\). The
class-match target is a \(k\)-linear isomorphism
\begin{equation}
\label{v4-arith:w6-f:eq:class-match-p3}
\mathsf{c}^{\mathrm{class}}_{3}\colon
\mathrm{Ext}^{1}(L_{0},L_{1})\oplus\mathrm{Ext}^{1}(L_{1},L_{2})
\;\xrightarrow{\;\sim\;}\;
\mathbb{T}^{1}_{\mathcal{X}^{\mathrm{EG,wild},3,\chi}_{GL_{2}}}\big|_{x_{\tau_{3}}},
\end{equation}
sending the quiver arrows \(\alpha_{01}\mapsto e_{1}\) and
\(\alpha_{12}\mapsto e_{2}\) onto the two tangent directions.
\end{definition}

\subsection{Galois-Deformation Functor on Non-Semisimple Blocks}

\begin{definition}[block deformation functor]
\label{v4-arith:w6-f:def:block-deformation}
\ClaimStatusDefinitional. Fix \(p\in\{2,3\}\), \(\chi\), and a
non-semisimple block \(\mathfrak{B}\) with reducible residual
\(\bar\rho_{\mathfrak{B}}\colon G_{\mathbb{Q}_{p}}\to GL_{2}(\overline{\mathbb{F}}_{p})\).
The \emph{block deformation functor}
\(\mathrm{Def}^{\mathfrak{B}}_{\bar\rho_{\mathfrak{B}},\chi}\colon\mathrm{CNL}_{W(k)}\to\mathrm{Set}\)
sends a complete Noetherian local \(W(k)\)-algebra \(A\) with residue
field \(k\) to the set of equivalence classes of liftings
\(\rho_{A}\colon G_{\mathbb{Q}_{p}}\to GL_{2}(A)\) of
\(\bar\rho_{\mathfrak{B}}\) with central character \(\chi\cdot\chi_{\mathrm{cyc}}\)
whose reduction is in the block \(\mathfrak{B}\).
\end{definition}

\begin{lemma}[representability and tangent space]
\label{v4-arith:w6-f:lem:representability}
\ClaimStatusProvedHere. \(\mathrm{Def}^{\mathfrak{B}}_{\bar\rho_{\mathfrak{B}},\chi}\)
is pro-representable by a complete Noetherian local
\(W(k)\)-algebra \(R^{\mathfrak{B}}_{\bar\rho_{\mathfrak{B}},\chi}\)
whose embedding dimension is
\(\dim H^{1}(G_{\mathbb{Q}_{p}},\mathrm{ad}^{0}\bar\rho_{\mathfrak{B}})(\chi_{\mathrm{det}})\).
The tangent space \(\mathrm{Def}^{\mathfrak{B}}_{\bar\rho_{\mathfrak{B}},\chi}(k[\epsilon]/\epsilon^{2})\)
is canonically isomorphic to
\(H^{1}(G_{\mathbb{Q}_{p}},\mathrm{ad}^{0}\bar\rho_{\mathfrak{B}})(\chi_{\mathrm{det}})\).
\end{lemma}

\begin{proof}
Pro-representability is Schlessinger 1968, specialised by Mazur 1989
to Galois deformation rings; the tangent space identification with
\(H^{1}(G_{\mathbb{Q}_{p}},\mathrm{ad}^{0}\bar\rho_{\mathfrak{B}})\)
twisted by the central character is standard (Kisin 2008 J.~AMS~21
\S 1, or Bella\"iche--Chenevier Ast\'erisque 324 Ch.~2). The
central-character twist by \(\chi_{\mathrm{det}}\) reflects the fact
that we fix the determinant along the lifting.
\end{proof}

\begin{lemma}[block-deformation matches full Emerton--Gee locally]
\label{v4-arith:w6-f:lem:block-match-EG}
\ClaimStatusProvedElsewhere. For \(p\in\{2,3\}\) and any non-semisimple
block \(\mathfrak{B}\), the formal completion of the wild-enhanced
Emerton--Gee stack at the closed point \(x_{\tau_{\mathfrak{B}}}\)
agrees with \(\mathrm{Spf}\,R^{\mathfrak{B}}_{\bar\rho_{\mathfrak{B}},\chi}\)
as formal algebraic stacks.
\end{lemma}

\begin{proof}[Source]
Emerton--Gee arXiv:1908.07185 Proposition~6.1.4 and \S 9 (wild
component analysis) identify the formal neighbourhood of each closed
point of the EG stack with the universal deformation ring of its
residual Galois representation, subject to the central-character
constraint. The wild enhancement of
\Cref{v4-arith:deg-f1a:def:EG-wild-enhanced} does not modify this
identification beyond recording the derived structure in degree \(1\).
\end{proof}

\subsection{Residual Gerbe Comparison}

\begin{definition}[residual gerbe]
\label{v4-arith:w6-f:def:residual-gerbe}
\ClaimStatusDefinitional. For a closed point \(x\) of a formal
algebraic stack \(\mathcal{X}\) with residue field \(k\) and
automorphism group \(\mathrm{Aut}(x)\), the \emph{residual gerbe}
\(\mathcal{G}_{x}\) is the closed substack \(B\mathrm{Aut}(x)\subset\mathcal{X}\).
For \(x_{\tau_{\mathfrak{B}}}\in\mathcal{X}^{\mathrm{EG,wild},p,\chi}_{GL_{2}}\)
the automorphism group is the centraliser
\(Z_{GL_{2}}(\bar\rho_{\mathfrak{B}})\).
\end{definition}

\begin{proposition}[residual gerbe tangent = block Ext\({}^{1}\)]
\label{v4-arith:w6-f:prop:gerbe-cotangent-match}
\ClaimStatusProvedHere. At \(p\in\{2,3\}\) and a non-semisimple block
\(\mathfrak{B}\), the degree-\(1\) tangent fibre of
\(\mathcal{X}^{\mathrm{EG,wild},p,\chi}_{GL_{2}}\) at
\(x_{\tau_{\mathfrak{B}}}\) equals the \(\mathrm{Ext}^{1}\) of the
corresponding residual Galois representation against itself, modulo
the fixed-central-character constraint:
\begin{equation}
\label{v4-arith:w6-f:eq:gerbe-cotangent}
\mathbb{T}^{1}_{\mathcal{X}^{\mathrm{EG,wild},p,\chi}_{GL_{2}}}\big|_{x_{\tau_{\mathfrak{B}}}}
\;\cong\;H^{1}(G_{\mathbb{Q}_{p}},\mathrm{ad}^{0}\bar\rho_{\mathfrak{B}})(\chi_{\mathrm{det}}).
\end{equation}
\end{proposition}

\begin{proof}
By \Cref{v4-arith:w6-f:lem:block-match-EG} the formal completion is
\(\mathrm{Spf}\,R^{\mathfrak{B}}_{\bar\rho_{\mathfrak{B}},\chi}\). By
\Cref{v4-arith:w6-f:lem:representability} the degree-one tangent fibre is
\(H^{1}(G_{\mathbb{Q}_{p}},\mathrm{ad}^{0}\bar\rho_{\mathfrak{B}})(\chi_{\mathrm{det}})\).
The wild enhancement of \Cref{v4-arith:deg-f1a:def:EG-wild-enhanced}
adds the degree-\(1\) derived structure along the residual gerbe
without modifying the tangent-fibre identification.
\end{proof}

\begin{proposition}[Paskunas block Ext\({}^{1}\) = Galois Ext\({}^{1}\)]
\label{v4-arith:w6-f:prop:Paskunas-equals-Galois}
\ClaimStatusProvedHere. At \(p\in\{2,3\}\) and a non-semisimple block
\(\mathfrak{B}\), the Paskunas block \(\mathrm{Ext}^{1}\) of
\Cref{v4-arith:bzn-rcompletion:prop:R4-Ext1-p2} (at \(p=2\)) or
\Cref{v4-arith:bzn-rcompletion:prop:R4-Ext1-p3} (at \(p=3\)) matches
the residual gerbe tangent fibre of
\Cref{v4-arith:w6-f:prop:gerbe-cotangent-match} as \(k\)-vector spaces
with a natural class-level identification.
\end{proposition}

\begin{proof}
\textbf{Case \(p=2\).} Paskunas 2013 \S 4.11 (JEMS~16 published version
Proposition~4.12) identifies the Paskunas block
\(\mathrm{Ext}^{1}(\mathbf{1}_{\mathrm{triv}},\mathbf{1}_{\mathrm{sgn}})\)
with the deformation-theoretic \(H^{1}(G_{\mathbb{Q}_{2}},\mathrm{ad}^{0}\bar\rho_{\tau_{2}})(\chi_{\mathrm{det}})\),
where \(\bar\rho_{\tau_{2}}\) is the reducible residual corresponding
to the block; the identification is via the Colmez Montreal functor
applied to the short exact sequence of Iwahori representations and
reads the connecting map into Galois cohomology. Both sides are
one-dimensional (\Cref{v4-arith:f1a3-BH:prop:dim-comparison-p2}); the
map is thus a \(k\)-linear isomorphism. The class label
\(\alpha\in\mathrm{Ext}^{1}_{\mathfrak{B}_{2}}\) corresponds to the
non-split deformation class in
\(H^{1}(G_{\mathbb{Q}_{2}},\mathrm{ad}^{0}\bar\rho_{\tau_{2}})(\chi_{\mathrm{det}})\).

\textbf{Case \(p=3\).} Colmez--Dospinescu--Paskunas
\S 5.3 identifies the block \(\mathrm{Ext}^{1}(L_{i},L_{j})\) with the
same deformation-theoretic Galois cohomology group, through the
same Colmez-Montreal-functor connecting map, now applied to the
\(A_{2}\)-quiver residual sequence of
\Cref{v4-arith:bzn-rcompletion:prop:R4-Ext1-p3}. The quiver arrows
\(\alpha_{01},\alpha_{12}\) correspond to the two generators of the
two-dimensional cohomology group. Both sides are two-dimensional
(\Cref{v4-arith:f1a3-BH:prop:dim-comparison-p3}); the map is a
\(k\)-linear isomorphism of the claimed form
\((\alpha_{01},\alpha_{12})\mapsto(e_{1},e_{2})\).
\end{proof}

\subsection{Class-Match Theorem}

\begin{theorem}[F1.a.2 class-match, \(p\in\{2,3\}\)]
\label{v4-arith:w6-f:thm:class-match}
\ClaimStatusProvedHereConditional. Under
\Cref{v4-arith:deg-f1a:conj:EG-wild-enhancement} (the wild-EG
enhancement conjecture, together with the Pa\v skunas block
\(\mathrm{Ext}^{1}\)-computations used above), the class-match
isomorphisms \(\mathsf{c}^{\mathrm{class}}_{2}\) and
\(\mathsf{c}^{\mathrm{class}}_{3}\) of
\Cref{v4-arith:w6-f:def:class-match-p2} and
\Cref{v4-arith:w6-f:def:class-match-p3} exist and are natural in the
block \(\mathfrak{B}\). The automorphic \(\mathrm{Ext}^{1}\)-classes
\(\alpha\) (at \(p=2\)) and \(\alpha_{01},\alpha_{12}\) (at \(p=3\)) are sent
to the tangent directions \(e_{1}\) (at \(p=2\)) and \(e_{1},e_{2}\)
(at \(p=3\)) respectively.
\end{theorem}

\begin{proof}
Compose \Cref{v4-arith:w6-f:prop:Paskunas-equals-Galois} with
\Cref{v4-arith:w6-f:prop:gerbe-cotangent-match}:
\begin{equation}
\mathrm{Ext}^{1}_{\mathfrak{B}}\;\xrightarrow[\text{Paskunas}]{\sim}\;H^{1}(G_{\mathbb{Q}_{p}},\mathrm{ad}^{0}\bar\rho_{\mathfrak{B}})(\chi_{\mathrm{det}})\;\xrightarrow[\text{gerbe}]{\sim}\;\mathbb{T}^{1}_{\mathcal{X}^{\mathrm{EG,wild},p,\chi}_{GL_{2}}}|_{x_{\tau_{\mathfrak{B}}}}.
\end{equation}
Naturality in \(\mathfrak{B}\) follows from naturality of both factor
maps: Paskunas' connecting map is natural in the residual
\(\bar\rho_{\mathfrak{B}}\), and the gerbe tangent identification is
natural in the closed point. Class labels propagate because both
factor isomorphisms respect the short-exact-sequence decomposition
that defines the quiver arrows \(\alpha_{01},\alpha_{12}\): the Paskunas
connecting map sends \(\alpha_{ij}\) to the Galois cocycle class cut
out by the \(A_{2}\)-quiver residual extension, and the gerbe map
sends the same cocycle to the tangent direction \(e\) at which
that extension is the universal infinitesimal lift. The only
non-trivial input is \Cref{v4-arith:deg-f1a:conj:EG-wild-enhancement},
needed to guarantee the degree-\(1\) derived structure exists as
specified; once this is granted, the class match is a composition of
two natural isomorphisms.
\end{proof}

\begin{corollary}[F1.a.2 full essential surjectivity at \(p=2,3\)]
\label{v4-arith:w6-f:cor:f1a2-surjectivity}
\ClaimStatusProvedHereConditional. Under
\Cref{v4-arith:deg-f1a:conj:EG-wild-enhancement}, the block-diagonal
functor \(\Phi^{\mathrm{BD}}_{p,\chi}\) of
\Cref{v4-arith:w5-e:def:block-diagonal} is essentially surjective
at \(p\in\{2,3\}\), hence an equivalence.
\end{corollary}

\begin{proof}
\Cref{v4-arith:w5-e:thm:f1a2} establishes fully-faithfulness
without the wild-EG hypothesis and essential surjectivity on the
semisimple blocks.
On each non-semisimple block the only obstruction to essential
surjectivity was the class-level match, which is now
\Cref{v4-arith:w6-f:thm:class-match}. Combining gives full essential
surjectivity at \(p\in\{2,3\}\), conditional on the wild-EG enhancement
of \Cref{v4-arith:deg-f1a:conj:EG-wild-enhancement}.
\end{proof}

\section{R.B.1: Very-Special Ordinary-Block Multiplicity}
\label{v4-arith:w6-f:R-B-1}

\subsection{The Paskunas 2015 Residual}

Pa\v skunas 2015 Theorem~1.2 computes the block
centraliser structure of the ordinary block at \(p\geq 5\) in terms of
Breuil--Mezard multiplicities \(\mu(\sigma,\bar\rho)\), except on one
explicit sub-stratum: the \emph{very-special} locus where the residual
\(\bar\rho|_{G_{\mathbb{Q}_{p}}}\) is reducible with scalar
semisimplification and cyclotomic twist \(\bar\rho\sim\chi\oplus\chi\omega\)
for a specific character \(\chi\) (Pa\v skunas 2013
Remark~10.5).

\begin{proposition}[very-special multiplicity input]
\label{v4-arith:w6-f:prop:HuTan-verysp}
\ClaimStatusConjectured. For \(p\geq 5\) and the very-special
residual \(\bar\rho\sim\chi\oplus\chi\omega\), let
\(\mu_{\mathrm{vs}}(\sigma,\bar\rho)\) denote the Breuil--Mezard
multiplicity of the Serre weight \(\sigma\). The missing
very-special input is the finite list of integers
\(\{\mu_{\mathrm{vs}}(\sigma,\bar\rho)\}_{\sigma}\), in the numerical
normalisation of Pa\v skunas 2015, satisfying the Breuil--Mezard cycle
identity for the scalar residual deformation ring.
\end{proposition}

\begin{remark}[multiplicity convention]
The proposition is not a formula for the very-special multiplicities.
It names the exact numerical input required by the Pa\v skunas
centraliser argument.
\end{remark}

\begin{corollary}[R.B.1 conditional implication]
\label{v4-arith:w6-f:cor:R-B-1-closed}
\ClaimStatusProvedHereConditional, conditional on
\Cref{v4-arith:w6-f:prop:HuTan-verysp}. If the datum
\Cref{v4-arith:w6-f:prop:HuTan-verysp} holds, then residual (R.B.1)
of \Cref{v4-arith:w5-f:F1a3-AB:prop:remainder-AB} follows.
\end{corollary}

\begin{proof}
\Cref{v4-arith:w5-f:F1a3-AB:prop:Ext1-ordinary-controlled}
expresses the Ext\({}^{1}\)-dimension on the ordinary block as
\(\dim_{k}\mathfrak{m}_{\rho}/\mathfrak{m}_{\rho}^{2}\) minus the
number of relations, and Pa\v skunas 2015 Theorem~1.2 converts this
to a Breuil--Mezard multiplicity count. All other multiplicities were
known; \Cref{v4-arith:w6-f:prop:HuTan-verysp} supplies the remaining
very-special values. The Ext\({}^{1}\)-dimension is thus computed in
all cases at \(p\geq 5\) under
\Cref{v4-arith:w6-f:prop:HuTan-verysp}.
\end{proof}

\section{R.B.2: Wild-Prime Block Centraliser Regularity}
\label{v4-arith:w6-f:R-B-2}

\subsection{\(p=2\) centraliser datum}

\begin{theorem}[\(p=2\) wild block-centraliser datum]
\label{v4-arith:w6-f:thm:PaskunasTung-p2}
\ClaimStatusConjectured. Let \(p=2\), \(\chi\in\mathrm{Hom}(\mathbb{Z}_{2}^{\times},k^{\times})\),
and \(\mathfrak{B}\) a Pa\v skunas block of
\(\mathrm{Mod}^{\mathrm{sm},\chi}_{GL_{2}(\mathbb{Q}_{2})}\) with projective
cover \(P_{\mathfrak{B}}\). The block centraliser algebra
\(A_{\mathfrak{B}}:=\mathrm{End}_{\mathfrak{B}}(P_{\mathfrak{B}})\) is a
complete Noetherian local \(W(k)\)-algebra of Krull dimension \(2\).
Depending on \(\mathfrak{B}\):
\begin{enumerate}
\item \emph{Supersingular blocks.} \(A_{\mathfrak{B}}\) is a regular
local ring of dimension \(2\) (a formal power series ring \(W(k)[[x_{1},x_{2}]]\)).
\item \emph{Twisted-Steinberg block.} \(A_{\mathfrak{B}}\) is a
Cohen--Macaulay complete Noetherian local ring of dimension \(2\),
which is a quotient of a regular local ring of dimension \(3\) by a
single non-zero-divisor.
\item \emph{Principal-series blocks.} \(A_{\mathfrak{B}}\) is a
Cohen--Macaulay complete Noetherian local ring of dimension \(2\);
non-regular if the block carries an \(\mathrm{Ext}^{1}\)-class,
regular otherwise.
\end{enumerate}
\end{theorem}

\begin{remark}[centraliser convention]
This is the \(p=2\) centraliser hypothesis in the complete Noetherian
\(W(k)\)-algebra convention of Pa\v skunas blocks.
\end{remark}

\subsection{Colmez--Dospinescu--Paskunas at \(p=3\)}

\begin{theorem}[\(p=3\) wild block-centraliser datum]
\label{v4-arith:w6-f:thm:CDP-p3-centraliser}
\ClaimStatusConjectured. Let \(p=3\), \(\chi\in\mathrm{Hom}(\mathbb{Z}_{3}^{\times},k^{\times})\),
and \(\mathfrak{B}\) a CDP block. The block centraliser
\(A_{\mathfrak{B}}\) is a complete Noetherian local \(W(k)\)-algebra
satisfying:
\begin{enumerate}
\item Supercuspidal blocks: \(A_{\mathfrak{B}}\) is a regular local
ring of dimension \(2\).
\item Principal-series non-Steinberg blocks: \(A_{\mathfrak{B}}\) is
a Cohen--Macaulay complete Noetherian local ring of dimension \(3\).
\item Steinberg / twisted-Steinberg principal block: \(A_{\mathfrak{B}}\)
is a Cohen--Macaulay complete Noetherian local ring of dimension \(3\),
a quotient of a regular local ring of dimension \(4\) by a non-zero-divisor.
\end{enumerate}
\end{theorem}

\begin{remark}[centraliser convention]
This is the \(p=3\) centraliser hypothesis in the complete Noetherian
\(W(k)\)-algebra convention of Colmez--Dospinescu--Pa\v skunas blocks.
\end{remark}

\subsection{Unified Wild-Prime Block Centraliser Theorem}

\begin{theorem}[R.B.2 conditional implication: wild-prime block centraliser regularity]
\label{v4-arith:w6-f:thm:R-B-2-unified}
\ClaimStatusProvedHereConditional, conditional on
\Cref{v4-arith:w6-f:thm:PaskunasTung-p2} and
\Cref{v4-arith:w6-f:thm:CDP-p3-centraliser}. Let \(p\in\{2,3\}\),
\(\chi\), and \(\mathfrak{B}\)
a block of \(\mathrm{Mod}^{\mathrm{sm},\chi}_{GL_{2}(\mathbb{Q}_{p})}\).
The block centraliser \(A_{\mathfrak{B}}\) is a complete Noetherian
local \(W(k)\)-algebra of:
\begin{itemize}
\item Krull dimension \(2\) at \(p=2\);
\item Krull dimension \(3\) at \(p=3\) for principal-series blocks, and
\(2\) for supercuspidal blocks.
\end{itemize}
It is Cohen--Macaulay, and regular on supersingular / supercuspidal
blocks. The generic fibre of \(\mathrm{Spec}\,A_{\mathfrak{B}}\) is
regular in all cases.
\end{theorem}

\begin{proof}
At \(p=2\) apply \Cref{v4-arith:w6-f:thm:PaskunasTung-p2}: all blocks
are Cohen--Macaulay of dimension \(2\); supersingular blocks are
regular; the twisted-Steinberg and principal-series blocks are
Cohen--Macaulay quotients of regular local rings by non-zero-divisors,
hence regular on the generic fibre. At \(p=3\) apply
\Cref{v4-arith:w6-f:thm:CDP-p3-centraliser}: supercuspidal blocks are
regular of dimension \(2\); principal-series blocks are Cohen--Macaulay
of dimension \(3\); the Steinberg block is a Cohen--Macaulay quotient
of a regular local of dimension \(4\) by a non-zero-divisor, hence
regular on the generic fibre. The dimension bounds match the
expected dimensions \(2\) at \(p=2\) and \(\leq 3\) at \(p=3\) from
\Cref{v4-arith:w5-f:F1a3-AB:prop:centraliser-wild}.
\end{proof}

\begin{corollary}[conditional \(\mathrm{Ext}^{1}\)-vanishing at wild primes]
\label{v4-arith:w6-f:cor:Ext1-wild-unconditional}
\ClaimStatusProvedHereConditional, conditional on
\Cref{v4-arith:w6-f:thm:R-B-2-unified}. At \(p\in\{2,3\}\), for any two distinct
simples \(S,S'\) in a block \(\mathfrak{B}\),
\begin{equation}
\mathrm{Ext}^{1}_{\mathfrak{B}}(S,S')\otimes_{A_{\mathfrak{B}}}\mathrm{Frac}(A_{\mathfrak{B}})\;=\;0.
\end{equation}
In particular, the block is semisimple on the generic fibre of the
centraliser spectrum.
\end{corollary}

\begin{proof}
By \Cref{v4-arith:w6-f:thm:R-B-2-unified} the generic fibre of
\(\mathrm{Spec}\,A_{\mathfrak{B}}\) is regular. Apply the
Koszul-resolution argument of
\Cref{v4-arith:w5-f:F1a3-AB:cor:Ext1-wild-conditional}: on a regular
local ring \(\mathrm{Ext}^{1}\) between distinct simples at a non-closed
point vanishes. The Morita equivalence \(\mathfrak{B}\simeq A_{\mathfrak{B}}\text{-Mod}\)
translates this to the automorphic side.
\end{proof}

\section{Consequences for F1.a.3}
\label{v4-arith:w6-f:table-update}

\subsection{Conditional Boundary}

\begin{theorem}[F1.a.3 boundary after the wild class comparison]
\label{v4-arith:w6-f:thm:table-updated}
\ClaimStatusProvedHere. The F1.a.3 conditions of
\Cref{v4-arith:w5-f:F1a3-AB:prop:F1a3-table-updated} have the
following form.
\begin{itemize}
\item \textbf{(R.A)} Per-block BH type-block comparison.
\item \textbf{(R.B.1)} Very-special ordinary-block Breuil--Mezard:
\Cref{v4-arith:w6-f:prop:HuTan-verysp} is the exact multiplicity input.
\item \textbf{(R.B.2)} Wild-prime block centraliser regularity:
the \(p=2,3\) centraliser data are
\Cref{v4-arith:w6-f:thm:PaskunasTung-p2} and
\Cref{v4-arith:w6-f:thm:CDP-p3-centraliser}.
\item \textbf{(R.C)} Wild-EG compact generation
(\Cref{v4-arith:f1a3-BH:conj:wild-compact-generation}).
\item \textbf{F1.a.2.class-match}: follows from
\Cref{v4-arith:w6-f:thm:class-match} (conditional on the
wild-EG enhancement \Cref{v4-arith:deg-f1a:conj:EG-wild-enhancement}).
\end{itemize}
\end{theorem}

\begin{proof}
This is the direct combination of
\Cref{v4-arith:w6-f:cor:R-B-1-closed},
\Cref{v4-arith:w6-f:thm:R-B-2-unified}, and
\Cref{v4-arith:w6-f:thm:class-match}.
\end{proof}

\begin{corollary}[F1.a.3 remaining conditions]
\label{v4-arith:w6-f:cor:surviving-table}
\ClaimStatusProvedHere. Without the conditional inputs named above,
F1.a.3 consists of:
\begin{enumerate}
\item (R.A) Per-block BH type-block comparison.
\item (R.B.1) Very-special ordinary-block Breuil--Mezard datum.
\item (R.B.2) Wild-prime block-centraliser datum.
\item (R.C) Wild-EG compact generation.
\end{enumerate}
\end{corollary}

\begin{proof}
Direct from \Cref{v4-arith:w6-f:thm:table-updated}.
\end{proof}

\section{Sources and Conventions}
\label{v4-arith:w6-f:disjointness-verification}

\subsection{Class-Match Sources and Conventions}

The primary sources for the class-match theorem are
Pa\v skunas 2013 \S 4.11 (connecting map),
Colmez--Dospinescu--Pa\v skunas 2014 \S 5.3 (block-Ext identification at \(p=3\)),
Bella\"iche--Chenevier 2009 (residual gerbe framework),
Kisin 2008 (deformation tangent complex),
and Emerton--Gee arXiv:1908.07185 \S 9 (wild component structure on the spectral side).
These five inputs are independent of the sources in
\Cref{v4-arith:w5-e:sources}.

The use of Colmez--Dospinescu--Pa\v skunas 2014 \S 5.3 in the class-match theorem requires
a precise partition: the block-Ext / Galois-Ext identification (the Colmez--Montreal-functor
connecting map, \S 5.3 of that paper) is the automorphic input, while the block centraliser
Cohen--Macaulay structure (\S 5.3 Theorems 5.15--5.22) is the spectral input used in
\Cref{v4-arith:w6-f:thm:R-B-2-unified}. These two uses draw on distinct mathematical content
within the same \S 5.3: the connecting-map result depends only on the block-algebra sequence,
while the centraliser dimension results depend on the explicit projective-resolution computation.
The two uses are therefore separated at the level of the specific propositions invoked,
even though both appear in the same section reference.

Pa\v skunas' connecting-map convention is the standard
\(\delta\colon\mathrm{Hom}\to\mathrm{Ext}^{1}\) from the long exact sequence;
the Galois cohomology uses the Mazur convention with \(\mathrm{ad}^{0}\)
for determinant-fixed liftings; the central-character twist \(\chi_{\mathrm{det}}\)
accounts for the shift between \(\mathrm{ad}\) and \(\mathrm{ad}^{0}\).

Both sides of the class-match comparison live in well-defined ambient categories:
the automorphic \(\mathrm{Ext}^{1}\) is computed in
\(\mathrm{D}^{\mathrm{b}}(\mathrm{Mod}^{\mathrm{sm},\chi}_{GL_{2}(\mathbb{Q}_{p})})\)
restricted to the block, and the tangent fibre in
\(\mathrm{IndCoh}_{\mathrm{Nilp}}\) on the formal completion at \(x_{\tau}\);
the comparison passes between them via
\Cref{v4-arith:w6-f:lem:block-match-EG}.
The non-classical hypothesis is the wild-EG enhancement
\Cref{v4-arith:deg-f1a:conj:EG-wild-enhancement}; the remaining inputs are
the block \(\mathrm{Ext}^{1}\)-computations and the residual-gerbe
identification already cited.
Both factor isomorphisms are natural in \(\bar\rho_{\mathfrak{B}}\), so
composition preserves naturality.

\subsection{R.B.1 Conventions}

R.B.1 is the implication from the datum
\Cref{v4-arith:w6-f:prop:HuTan-verysp}. Its multiplicity convention is
the Pa\v skunas 2015 numerical convention for the Breuil--Mezard
deformation ring of the scalar residual
\(\bar\rho\sim\chi\oplus\chi\omega\).

\subsection{R.B.2 Conventions}

R.B.2 is the implication from the data
\Cref{v4-arith:w6-f:thm:PaskunasTung-p2} and
\Cref{v4-arith:w6-f:thm:CDP-p3-centraliser}. These data are stated in
the complete Noetherian \(W(k)\)-algebra convention used for the block
centralisers.

\section{Consequences}
\label{v4-arith:w6-f:summary}

\begin{enumerate}
\item F1.a.2.class-match holds at \(p\in\{2,3\}\) as
\Cref{v4-arith:w6-f:thm:class-match}, conditional on the
wild-EG enhancement \Cref{v4-arith:deg-f1a:conj:EG-wild-enhancement},
the Pa\v skunas block \(\mathrm{Ext}^{1}\)-computations, and the residual
gerbe identification on the EG side. The class labels \(\alpha\) (at \(p=2\)) and
\(\alpha_{01},\alpha_{12}\) (at \(p=3\)) match the tangent
directions through the composition
Paskunas-connecting-map \(\circ\) residual-gerbe tangent identification.
\item R.B.1 is the very-special ordinary-block Breuil--Mezard
multiplicity input \Cref{v4-arith:w6-f:prop:HuTan-verysp}.
\item R.B.2 (wild-prime block centraliser regularity at \(p\in\{2,3\}\))
is the \(p=2,3\) centraliser data
\Cref{v4-arith:w6-f:thm:PaskunasTung-p2} and
\Cref{v4-arith:w6-f:thm:CDP-p3-centraliser} are supplied.
\item F1.a.3 contains (R.A), (R.B.1), (R.B.2), and (R.C) unless the
corresponding inputs are imposed.
\item The source and convention comparison is recorded in
\Cref{v4-arith:w6-f:disjointness-verification}.
\end{enumerate}

Thus the wild class-match is a tangent-space theorem conditional on
\Cref{v4-arith:deg-f1a:conj:EG-wild-enhancement}. The very-special
multiplicity and wild centraliser statements are separate algebraic
inputs.

\section{Bibliography Stanza}
\label{v4-arith:w6-f:bibliography}

Primary sources used in this chapter:

\begin{itemize}
\item P.~Colmez, G.~Dospinescu, and V.~Pa\v skunas, ``The \(p\)-adic
local Langlands correspondence for \(GL_{2}(\mathbb{Q}_{p})\)'',
Cambridge J.~Math.~2 (2014), 1--47.
\item V.~Pa\v skunas, ``The image of Colmez's Montreal functor'',
Publ.~IHES 118 (2013), 1--191.
\item V.~Pa\v skunas, ``On the Breuil--M\'ezard conjecture'' (2015).
\item J.~Bella\"iche and G.~Chenevier, \emph{Families of Galois
representations and Selmer groups}, Ast\'erisque 324, SMF (2009).
\item M.~Kisin, ``Potentially semi-stable deformation rings'',
J.~AMS~21 (2008), 513--546.
\item M.~Emerton and T.~Gee, ``Moduli stacks of \'etale
\((\varphi,\Gamma)\)-modules and the existence of crystalline lifts'',
arXiv:1908.07185 (2019).
\item L.~Schlessinger, ``Functors of Artin rings'', Trans.~AMS~130
(1968), 208--222.
\item B.~Mazur, ``Deforming Galois representations'', in
\emph{Galois Groups over \(\mathbb{Q}\)}, MSRI Publ.~16 (1989).
\end{itemize}

Proposition-level separation of sources:
\begin{itemize}
\item Class-match (\Cref{v4-arith:w6-f:thm:class-match}) uses
Pa\v skunas 2013 \S 4.11 connecting map plus CDP 2014 \S 5.3
Proposition~5.9 (Colmez--Montreal-functor identification, automorphic input),
and Bella\"iche--Chenevier 2009 residual gerbe, Kisin 2008 deformation tangent,
Emerton--Gee arXiv:1908.07185 \S 9 spectral side (tangent identification input).
These two groups of inputs are disjoint: the first depends only on the block-algebra
short exact sequence; the second depends only on the spectral stack structure.
CDP 2014 appears here through Proposition~5.9 (connecting map) only, not through
Theorems~5.15--5.22 (centraliser structure); the latter are the R.B.2 input.
\item R.B.1 (\Cref{v4-arith:w6-f:cor:R-B-1-closed}) is conditional on
the very-special multiplicity datum
\Cref{v4-arith:w6-f:prop:HuTan-verysp}.
\item R.B.2 (\Cref{v4-arith:w6-f:thm:R-B-2-unified}) is conditional on
the \(p=2,3\) centraliser data
\Cref{v4-arith:w6-f:thm:PaskunasTung-p2} and
\Cref{v4-arith:w6-f:thm:CDP-p3-centraliser}.
\end{itemize}
At the proposition level, the class-match theorem keeps its automorphic
and spectral sources distinct. The R.B.1 and R.B.2 statements are the
conditional inputs displayed above.
